\documentclass[10pt,a4paper]{amsart}
\usepackage[margin=19mm]{geometry}
\usepackage{amsmath,amssymb,mathtools}
\usepackage[scr=rsfs,cal=euler]{mathalfa}
\usepackage{xfrac}
\usepackage[unicode,hidelinks,bookmarks=false]{hyperref}
\newcommand{\ie}{i.e.\ }
\newcommand{\cf}{cf.\ }
\newcommand{\resp}{respectively\ }
\newcommand{\Subsection}{\subsection}
\providecommand{\nobreakdash}{\nobreak\hbox{-}}
\setlength{\emergencystretch}{2em}
\multlinegap0pt
\newtheorem{theorem}{Theorem}
\newtheorem{lemma}{Lemma}
\newtheorem{corollary}{Corollary}
\newcommand{\Psu}{\mathrm{P}_{\mathcal{U}}}
\newcommand{\Fit}{\mathrm{Fit}}
\newcommand{\cp}{\mathrm{cp}}
\newcommand{\Frat}{\Phi}
\begin{document}
\title[Supersoluble groups and the probability of generating a supe]{Supersoluble groups and the probability of generating a supersoluble subgroup}
\author{Andrea Lucchini Andrea Lucchini; Dipartimento di Matematica Tullio Levi-Civitaq̊q̊; Andrea Lucchini}
\date{}
\maketitle
\noindent\emph{Source and licence.} Supersoluble groups and the probability of generating a supersoluble subgroup. Version arXiv:2607.08165v1 (CC BY 4.0). This material is distributed under the Creative Commons Attribution 4.0 International licence, http://creativecommons.org/licenses/by/4.0/, as recorded in the arXiv OAI licence metadata for this exact version and shown on the abstract page https://arxiv.org/abs/2607.08165v1. Full rights notice: licenses/arxiv-2607.08165-CC-BY-4.0.txt.\par\smallskip
\noindent\textit{Editorial scope.} Counted unit: the original \texttt{theorem} environment \texttt{thm:main} at source lines 80--87 together with its complete proof at lines 384--450; the delivered document keeps source lines 35--451 so that every referenced label and every citation resolves inside it. Native editable source is the paper's own concatenated TeX; the AMS wrapper is editorial formatting only. No publisher class, upstream script or unrelated artwork is redistributed.
	\section{Introduction}
	
	Let $\mathcal{C}$ be a class of finite groups. We say that 
	$\mathcal{C}$ is \textit{recognizable} if, whenever all 
	$2$-generated subgroups of a finite group $G$ belong to 
	$\mathcal{C}$, then $G$ itself belongs to $\mathcal{C}$. 
	It is well known for example that the class of finite soluble groups is recognizable (see Thompson \cite{th} and Flavell \cite{pf}). On the other hand the class of finite metabelian groups is not recognizable (see \cite[Example 7.1]{bln}).
	More quantitatively, for $0 < \eta < 1$, we say that 
	$\mathcal{C}$ is \textit{$\eta$-recognizable} if, whenever 
	the probability that a $2$-generated subgroup of $G$ belongs 
	to $\mathcal{C}$ is greater than $\eta$, then $G \in \mathcal{C}$. 
	Formally, letting
	$$\mathrm{P}_{\mathcal{C}}(G) = \frac{|\{(x,y) \in G \times G 
		\mid \langle x,y\rangle \in \mathcal{C}\}|}{|G|^2},$$
	we say that $\mathcal{C}$ is $\eta$-recognizable if 
	$\mathrm{P}_{\mathcal{C}}(G) > \eta$ implies $G \in \mathcal{C}$.
	
	Clearly, every $\eta$-recognizable class is recognizable. 
	It is natural to ask whether every recognizable class 
	is $\eta$-recognizable for some $\eta$, and what the optimal 
	value of $\eta$ is.
	
	Guralnick and Wilson \cite{GW00} proved that for any class 
	$\mathcal{C}$ closed under taking subgroups, quotients and 
	extensions, $\mathcal{C}$ is $\eta$-recognizable for some 
	$\eta > 0$ (combining \cite[Proposition 5]{mqrd} with \cite[Theorem
	1.1]{mqrd}, one may deduce that $\eta$ can be taken to be $\frac{37}{90})$.
	In particular they proved that the class 
	of finite soluble groups is $\frac{11}{30}$-recognizable
	and that this bound is sharp. 
	
	Earlier, 
	Gustafson \cite{Gus73} had proved that the class of finite  abelian 
	groups is $\frac{5}{8}$-recognizable, while Guralnick and Wilson in \cite{GW00} also prove that the class of finite nilpotent groups is $\frac{1}{2}$-recognizable. The next natural candidate to investigate is the class 
	$\mathcal{U}$ of supersoluble groups, which is not closed under taking extension and lies strictly 
	between the classes of nilpotent and soluble groups. 
	It is known that $\mathcal{U}$ is recognizable 
	(see for instance \cite[Example 1]{BH09}).
	 Our main result determines the 
	optimal threshold for $\eta$-recognizability.
	

	

	

	\begin{theorem}\label{thm:main}
		The class of supersoluble groups is $\frac{16}{25}$-recognizable. 
		More precisely, if $G$ is a finite group with 
		$\Psu(G) \geq \frac{16}{25}$, then $G$ is supersoluble. 
		Moreover, the bound $\frac{16}{25}$ is sharp, being attained 
		by the group $G = (C_5 \times C_5) \rtimes Q_8$, where $Q_8$ 
		acts faithfully and irreducibly on $C_5 \times C_5$.
	\end{theorem}

	
We note that it remains an open problem whether every 
recognizable class is $\eta$-recognizable for some 
$\eta > 0$. Natural candidates for further investigation 
are the class of groups with nilpotent derived subgroup 
and, more generally, the class of groups of Fitting length 
at most $k$, for a fixed positive integer $k$. These classes 
are known to be recognizable \cite[Examples 2 and 3]{BH09}, 
but whether they are $\eta$-recognizable for some $\eta > 0$ 
remains open.
	
	The results of Guralnick and Wilson easily reduce the proof of Theorem \ref{thm:main} to 
	 a detailed analysis of primitive soluble
	  groups with supersoluble point-stabilizers. 
	Recall that a finite group is \textit{primitive} if it has a 
	faithful primitive permutation representation; equivalently, 
	it has a core-free maximal subgroup. A primitive soluble group 
	has the form $G = V \rtimes H$, where $V = \mathrm{Soc}(G)$ is 
	elementary abelian and $H$ acts faithfully and irreducibly on $V$.
	
	The paper is organized as follows. In Section~2 we establish 
	the key lemmas on the module-theoretic structure of primitive 
	soluble groups. In Section~3 we prove the main theorem.
	
	\section{Preliminary lemmas}
	
	Throughout this section, $G = V \rtimes H$ is a primitive 
	soluble group, where $V$ is a faithful irreducible $H$-module 
	of order $p^n$, with $p$ prime and $n \geq 2$. We assume 
	moreover that $H$ is supersoluble.
	
	Since $H$ is supersoluble, $H/\Fit(H)$ is abelian. Since $V$ 
	is a faithful irreducible $H$-module, $O_p(H)=1$. Since 
	$P \cap \Fit(H) = 1$, a Sylow $p$-subgroup $P$ of $H$ is 
	isomorphic to its image in $H/\Fit(H)$, which is abelian, 
	hence $P$ is abelian. Moreover $H/\Fit(H)$ splits as
	$$H/\Fit(H) = K/\Fit(H) \times P\,\Fit(H)/\Fit(H)$$
	for a suitable subgroup $K$ of $H$, and $H = K \rtimes P$.
	
	For every pair $(h_1,h_2) \in H \times H$, let
	$$\Sigma(h_1,h_2) = \{(v_1,v_2)\in V^2 \mid 
	\langle v_1h_1, v_2h_2\rangle \text{ is supersoluble}\}$$
	and set $\pi(h_1,h_2) = |\Sigma(h_1,h_2)|/|V|^2$. Note that
	$$\Psu(G) = \frac{1}{|H|^2} \sum_{(h_1,h_2) \in H \times H} 
	\pi(h_1,h_2).$$
	
	We say that a pair $(h_1,h_2) \in H \times H$ is 
	\textit{$V$-large} if some composition factor of the 
	$\langle h_1,h_2\rangle$-module $V$ has order $\geq p^2$.
	
	\begin{lemma}\label{lem:large}
		If $(h_1,h_2)$ is $V$-large, then $\pi(h_1,h_2) \leq 1/p^2$.
	\end{lemma}
	\begin{proof}
		Let $X = \langle h_1,h_2\rangle$ and let $W_1 \leq W_2$ be 
		$X$-submodules of $V$ such that $W_2/W_1$ is $X$-irreducible 
		and $|W_2/W_1| \geq p^2$.
		
		Let $(v_1,v_2) \in V^2$ and set 
		$Y = \langle v_1h_1, v_2h_2\rangle$. Working modulo $W_1$, 
		if $Y$ is supersoluble then $YW_1/W_1$ is a complement of 
		$W_2/W_1$ in $\langle W_2/W_1, YW_1/W_1\rangle$. By a theorem 
		of Gasch\"{u}tz, all complements of $W_2/W_1$ are conjugate (see for example \cite[Propositon 1.6.8]{dh}), 
		so their number is $|W_2/W_1|$. Moreover, if 
		$\langle v_1h_1, v_2h_2\rangle W_1 = 
		\langle v_1'h_1, v_2'h_2\rangle W_1$ then 
		$(v_1-v_1', v_2-v_2') \in W_1 \times W_1$, so each complement 
		is realized by exactly $|W_1|^2$ pairs $(v_1,v_2) \in V^2$. 
		Therefore the number of pairs $(v_1,v_2)$ for which $Y$ is 
		supersoluble is at most 
		$$|W_2/W_1| \cdot |W_1|^2 = |W_2| \cdot |W_1| \leq 
		\frac{|W_2|^2}{p^2} \leq \frac{|V|^2}{p^2},$$
		where we used $|W_1| \leq |W_2|/p^2$ and $|W_2| \leq |V|$. 
		It follows that $\pi(h_1,h_2) \leq p^{-2}$.
	\end{proof}
	
	\begin{lemma}\label{lem:abelian}
		Let $X=\langle h_1,h_2\rangle \leq H$. If $X$ is not $V$-large, then 
		$X$ is abelian.
	\end{lemma}
	\begin{proof}
		If $(h_1,h_2)$ is not $V$-large, all the $X$-composition 
		factors of $V$ have order $p$. Let $\mathbb{F}$ be the field 
		of order $p$. Then there exists an $\mathbb{F}$-basis of $V$ 
		such that $\langle e_n,\dots,e_i\rangle$ is $X$-invariant for 
		every $1 \leq i \leq n$. With respect to this basis, $X$ 
		embeds into the group $T(n,p)$ of upper triangular $n\times n$ 
		matrices over $\mathbb{F}$. This implies that the Sylow 
		$p$-subgroup $Q$ of $X$ is normal in $X$ and has an abelian 
		complement $Y$. Since $X \leq K \rtimes P$, we have $Y \leq K$ 
		and $Q$ is abelian. Moreover 
		$[Q,Y] \leq [Q,K] \cap Q \leq K \cap Q = 1$, hence $X$ 
		is abelian.
	\end{proof}
	
	\begin{corollary}\label{cor:nonabelian}
		If $\langle h_1,h_2\rangle$ is nonabelian, then 
		$(h_1,h_2)$ is $V$-large, and in particular 
		$\pi(h_1,h_2) \leq 1/p^2$.
	\end{corollary}
	\begin{proof}
		Immediate from Lemmas~\ref{lem:abelian} and~\ref{lem:large}.
	\end{proof}
	
	
	\begin{lemma}\label{lem:q-large}
		Let $q$ be a prime power that does not divide $p-1$. If $q$ 
		divides the exponent $\exp(X)$ of $X$, then $(h_1, h_2)$ is $V$-large.
	\end{lemma}
	\begin{proof}
		Let $V = W_0 \geq W_1 \geq \dots \geq W_s = 1$ be a composition 
		series of $V$ as an $X$-module. The action of $X$ on each factor 
		$W_i/W_{i+1}$ induces a homomorphism
		$$\gamma \colon X \to \prod_{0 \leq i \leq s-1} \mathrm{Aut}(W_i/W_{i+1}),$$
		where the $i$-th component of $\gamma(x)$ is the automorphism 
		$wW_{i+1} \mapsto xwx^{-1}W_{i+1}$ of $W_i/W_{i+1}$.
		The kernel of $\gamma$ acts trivially on each factor 
		$W_i/W_{i+1}$, hence embeds in the group of upper triangular 
		matrices with ones on the diagonal over $\mathbb{F}_p$, which 
		is a $p$-group. Since $q$ is a prime power coprime to $p$, 
		$q$ divides $|\mathrm{Im}(\gamma)|$, hence $q$ divides 
		$|\mathrm{Aut}(W_i/W_{i+1})|$ for some $i$. Since 
		$|\mathrm{Aut}(W_i/W_{i+1})|$ divides $p-1$ when 
		$|W_i/W_{i+1}| = p$, and $q$ does not divide $p-1$, we 
		conclude that $|W_i/W_{i+1}| \geq p^2$, so $(h_1,h_2)$ 
		is $V$-large.
	\end{proof}
	
\begin{corollary}\label{cor:q-large}
	Let $q$ be a prime power that does not divide $p-1$. If $q$ 
	divides $\exp(\langle h_1, h_2 \rangle)$, then 
	$\pi(h_1,h_2) \leq 1/p^2$.
\end{corollary}
\begin{proof}
	Immediate from Lemmas~\ref{lem:q-large} and~\ref{lem:large}.
	\end{proof}


	\section{Proof of the main theorem}
	
	\subsection{The primitive case}
	
	We first establish the bound for primitive soluble groups.
	
	\begin{lemma}\label{lem:primitive}
		Let $G = V \rtimes H$ be a primitive soluble group that is 
		not supersoluble, where $|V| = p^n$ with $p$ prime and 
		$n \geq 2$, and $H$ is supersoluble. Then 
		$\Psu(G) \leq 16/25$, with equality if and only if $p = 5$, 
		$n=2$, and $H/Z(H) \cong C_2 \times C_2$.
	\end{lemma}
	\begin{proof}
		We consider separately the cases $H$ abelian and $H$ nonabelian.
		
		\medskip
		
		\noindent\textit{Case 1: $H$ abelian.}
		Since $H$ is abelian and acts faithfully and irreducibly on $V$, 
		$H$ embeds into $\mathrm{GL}(1,p^n) \cong \mathbb{F}_{p^n}^*$, 
		which is cyclic. Hence $H$ is cyclic of some order $m$ dividing 
		$p^n - 1$.
		
		Since $V$ is a faithful irreducible $H$-module and $n \geq 2$, 
		there exists a prime $q$ dividing $m$ such that $q$ divides 
		$p^n-1$ but does not divide $p-1$. In particular $p \neq q$.
		
		Write $m = q^a \cdot r$ with $\gcd(r,q) = 1$ and $a \geq 1$. 
		Since $H$ is cyclic of order $m$, the elements $h \in H$ whose 
		order is not divisible by $q$ are exactly those belonging to 
		the unique subgroup of order $r$. Therefore the probability 
		that $q$ does not divide $|\langle h_1,h_2\rangle|$ is at most
		$$\frac{r^2}{m^2} = \frac{1}{q^{2a}} \leq \frac{1}{q^2}.$$
		If $q$ divides $|\langle h_1,h_2\rangle|$, then %since $\langle h_1,h_2\rangle$ is cyclic (being a subgroup of the cyclic group $H$), $q$ also divides 	$\exp(\langle h_1,h_2\rangle)$, and so 
		$\pi(h_1,h_2) \leq 1/p^2$ by Corollary~\ref{cor:q-large}. 
		Therefore:
		\begin{align*}
			\Psu(G) &\leq \frac{1}{q^{2a}} + 
			\left(1-\frac{1}{q^{2a}}\right)\frac{1}{p^2}\\
			&= \frac{1}{p^2} + \frac{1}{q^{2a}}\left(1-\frac{1}{p^2}\right)\\
			&\leq \frac{1}{p^2}+\frac{1}{q^2}\left(1-\frac{1}{p^2}\right)\\
			&= \frac{1}{p^2}+\frac{1}{q^2}-\frac{1}{p^2q^2}\\
			&\leq \frac{1}{4}+\frac{1}{9}-\frac{1}{36}
			= \frac{1}{3} < \frac{16}{25},
		\end{align*}
		since $p$ and $q$ are distinct primes.
		
		
	
	\medskip
	\noindent\textit{Case 2: $H$ nonabelian.}	
		\medskip
		\noindent\textit{Case 2a: $p \geq 5$.}
			Let $\alpha$ denote the probability that two randomly chosen 
		elements of $H$ generate a nonabelian subgroup. By 
		Corollary~1, we get
		$$\mathrm{P}_{\mathcal{U}}(G) \leq \alpha \cdot \frac{1}{p^2} 
		+ (1-\alpha) \cdot 1 = 1 - \alpha\left(1-\frac{1}{p^2}\right).$$
	Let $\cp(H)$ denote the probability that two randomly chosen 
	elements of $H$ commute. By a theorem of Gustafson \cite{Gus73},
	$$\cp(H) \leq \frac{1}{2}\left(1 + \frac{|Z(H)|}{|H|}\right).$$
	Since $H$ is nonabelian, $|Z(H)|/|H|$ is not the reciprocal 
	of a prime, hence $|Z(H)|/|H| \leq 1/4$, which gives 
	$\cp(H) \leq 5/8$. Since every pair generating an abelian 
	subgroup commutes, $1 - \alpha \leq \cp(H)$, hence 
	$\alpha \geq 3/8$. Therefore,
	$$\Psu(G) \leq 1 - \frac{3}{8}\left(1-\frac{1}{p^2}\right) 
	\leq 1 - \frac{3}{8}\cdot\frac{24}{25} = \frac{16}{25},$$
	with equality if and only if $p=5$, $\alpha = 3/8$, and 
	$|Z(H)|/|H| = 1/4$, i.e.\ $H/Z(H) \cong C_2 \times C_2$, 
	which is realized by $H = Q_8$ and $H = D_4$.
		
		\medskip
		
		
		
		
		
		\noindent\textit{Case 2b: $p = 2$.}
	Since $V$ is a faithful irreducible $H$-module, $O_2(H)=1.$
	Hence $H/Z(H)$ is not a $2$-group, otherwise a Sylow 
	$2$-subgroup of $H$ would be normal, contradicting $O_2(H)=1$. 
	In particular $|H/Z(H)| \geq 6$ and $H$ contains an abelian subgroup 
	$K$ such that $Z(H) < K$ and $|K/Z(H)| \geq 3$.
	
	Since $p-1=1$, every prime $q$ satisfies $q \nmid p-1$, so 
	by Corollaries~\ref{cor:nonabelian} and~\ref{cor:q-large}, 
	$\pi(h_1,h_2) \leq 1/4$ if either $\langle h_1,h_2\rangle$ 
	is nonabelian or $1 \neq \langle h_1,h_2\rangle \leq K$.
	Setting $a = |H:Z(H)|$, it follows that
	\begin{align*}
		\Psu(G) &\leq 1-\frac{3}{4}\left(\left(1-\frac{1}{2}
		-\frac{1}{2a}\right)+\frac{|K|^2-1}{|H|^2}\right)\\
		&\leq 1-\frac{3}{4}\left(\frac{1}{2}-\frac{1}{2a}+
		\frac{9|Z(H)|^2-1}{|H|^2}\right)\\
		&\leq 1-\frac{3}{4}\left(\frac{1}{2}-\frac{1}{2a}+
		\frac{8}{a^2}\right)\\
		&\leq 1-\frac{3}{4}\left(\frac{1}{2}-\frac{1}{12}+
		\frac{8}{36}\right) = \frac{25}{48} < \frac{16}{25}.
	\end{align*}
	
		\medskip
\noindent\textit{Case 2c: $p = 3$.}
If $|H:Z(H)| \geq 6$, then
$$\Psu(G) \leq 1-\frac{8}{9}\left(1-\frac{1}{2}
-\frac{1}{2a}\right) \leq \frac{17}{27} < \frac{16}{25}.$$
So we may assume $H/Z(H) \cong C_2 \times C_2$. Since $Z(H)$ 
is cyclic (otherwise, by \cite[Proposition 16.8]{dh}, $H$ cannot 
have a faithful irreducible representation), we have 
$H = C \rtimes Q$ where $C$ is the unique subgroup of $Z(H)$ 
of odd order $m$ and $Q$ is a Sylow $2$-subgroup of $H$. 
Since $C \leq Z(H)$, if $Q$ were abelian then $H = C \times Q$ 
would be abelian, a contradiction. Hence $Q$ is nonabelian, 
with $Z(Q)$ cyclic of order $2^b$, $b \geq 1$.



Let $U$ be the unique subgroup of $Z(Q)$ of order $2$, and let $\iota$ denote the number of pairs 
$(h_1,h_2) \in Q \times Q$ with $\exp(\langle h_1,h_2\rangle) 
\leq 2$. Note that if $\exp(\langle h_1,h_2\rangle) \leq 2$ 
then $\langle h_1,h_2\rangle$ is abelian. If 
$\exp(\langle h_1,h_2\rangle) > 2$, then $\exp(\langle h_1,h_2\rangle)$ 
is divisible by some prime power $q \neq 2$ not dividing 
$p-1=2$, so by Corollary~\ref{cor:q-large}, 
$\pi(h_1,h_2)\leq 1/9$. Therefore
$$\Psu(G) \leq 1 - \frac{8}{9}\left(1 - \frac{\iota}{|H|^2}\right)
\leq 1 - \frac{8}{9}\left(1 - \frac{\iota}{|Q|^2}\right).$$
The pairs counted by $\iota$ belong neither to the at least 
$3|Q|^2/8$ pairs generating a nonabelian subgroup of $Q$, nor to 
the $2^{2b}-4$ pairs in $(Z(Q) \times Z(Q)) \setminus (U \times U$), 
hence)
$$\frac{\iota}{|Q|^2} \leq \frac{5}{8} - 
\frac{2^{2b}-4}{16\cdot 2^{2b}}.$$
For $b = 1$, $Q \cong Q_8$ or $Q \cong D_4$, and a direct 
computation gives $\iota/|Q|^2 = 1/16$ for $Q = Q_8$ 
and $\iota/|Q|^2 = 7/16$ for $Q = D_4$.
For $b = 2$, $Q \cong M_{16}$ or $Q \cong D_4 \ast C_4$, 
where $M_{16}$ denotes the modular maximal-cyclic group of 
order $16$ and $D_4 \ast C_4$ denotes the central product 
of $D_4$ with $C_4$. A direct computation gives 
$\iota/|Q|^2 = 1/16$ for $Q = M_{16}$ and 
$\iota/|Q|^2 = 5/32$ for $Q = D_4 \ast C_4$.
In all these cases $\iota/|Q|^2 \leq 7/16$, and therefore
$$\Psu(G) \leq 1 - \frac{8}{9}\cdot\frac{9}{16} = \frac{1}{2} 
< \frac{16}{25}.$$
For $b \geq 3$, using $b=3$ to minimize:
$$\Psu(G) \leq 1 - \frac{8}{9}\left(\frac{3}{8} + 
\frac{60}{1024}\right) = \frac{59}{96} < \frac{16}{25}.$$
	\end{proof}

		
	
	
	\subsection{Reduction to the primitive case}
	
	We now prove Theorem~\ref{thm:main} by induction on $|G|$.
	

	\begin{proof}[Proof of Theorem~\ref{thm:main}]
		Let $G$ be a finite group with $\Psu(G) \geq 16/25$. We 
		proceed by induction on $|G|$.
		
		\medskip
		\noindent\textit{Step 1: All proper quotients of $G$ are 
			supersoluble.}
		Let $N$ be a non-trivial normal subgroup of $G$. Since every 
		$2$-generated subgroup of $G/N$ is the image of a $2$-generated 
		subgroup of $G$, we have $\Psu(G/N) \geq \Psu(G) \geq 16/25$. 
		By induction, $G/N$ is supersoluble.
		
		\medskip
		\noindent\textit{Step 2: $G$ has a unique minimal normal 
			subgroup.}
		Suppose $G$ has two distinct minimal normal subgroups $N_1$ 
		and $N_2$. Then $N_1 \cap N_2 = 1$, so $G$ embeds into 
		$G/N_1 \times G/N_2$. Since both $G/N_1$ and $G/N_2$ are 
		supersoluble by Step~1, and the class of supersoluble groups 
		is closed under taking subgroups and direct products, $G$ is 
		supersoluble. 
		
		\medskip
		\noindent\textit{Step 3: $G$ is soluble.}
	The probability that two random elements of $G$ generate a 
	soluble subgroup is at least $\Psu(G) \geq 16/25 > 11/30$. 
	By a theorem of Guralnick and Wilson \cite[Theorem A]{GW00}, 
	the class of finite soluble groups is $\frac{11}{30}$-recognizable, 
	hence $G$ is soluble.
		
		\medskip
		\noindent\textit{Step 4: $G$ is primitive soluble.}
		By Steps~1 and~3, $G$ is soluble with a unique minimal normal 
		subgroup $N$. If $N$ is contained in the Frattini subgroup $\Phi(G)$ of $G$, then $G/\Phi(G)$ is a 
		quotient of $G/N$, hence supersoluble. Since supersolubility 
		is inherited from $G/\Phi(G)$ to $G$ \cite[VI Satz 8.6]{hup},
		$G$ would be supersoluble, a contradiction. Therefore 
		$N \not\leq \Phi(G)$, so there exists a maximal subgroup $M$ 
		of $G$ with $NM = G$ and $N \cap M = 1$. Since $N$ is the 
		unique minimal normal subgroup of $G$ and $\mathrm{Core}_G(M) 
		\leq N \cap M = 1$, the subgroup $M$ is core-free, and $G$ 
		is a primitive soluble group with $\mathrm{Soc}(G) = N$.
		
		\medskip
		\noindent\textit{Step 5: Conclusion.}
		Write $G = V \rtimes H$ where $V = N$ is elementary abelian 
		of order $p^n$ and $H \cong M$ is the point stabilizer. 
		Since $G/N \cong H$ is supersoluble by Step~1, and 
		$\Psu(G) \geq 16/25$, Lemma~\ref{lem:primitive} implies 
		that $G$ is supersoluble, completing the induction.
		
		\medskip
		\noindent\textit{Sharpness.}
The group $G = (C_5 \times C_5) \rtimes Q_8$, where $Q_8$ 
acts faithfully and irreducibly on $C_5 \times C_5$, is not 
supersoluble, since its unique minimal normal subgroup 
$C_5 \times C_5$ has order $25 = 5^2$, which is not prime. 
However, all proper subgroups of $G$ are supersoluble: indeed, 
since $5 \equiv 1 \pmod{4}$, the field $\mathbb{F}_5$ contains 
a primitive fourth root of unity, so every proper subgroup of 
$Q_8$ acts reducibly on $C_5 \times C_5$, and hence every 
proper subgroup of $G$ is supersoluble.
Since the probability that an ordered pair of elements 
generates $G$ is $9/25$, it follows that
$$\Psu(G) = 1 - 9/25 = 16/25.$$
This shows that the bound $16/25$ is sharp.
	\end{proof}
	

\begin{thebibliography}{99}
\bibitem[BH09]{BH09} J.~C.~Beidleman and H.~Heineken, \textit{Minimal non-$\mathfrak{F}$-groups}, Ric.\ Mat.\ \textbf{58} (2009), no.~1, 33--41.
\bibitem[GW00]{GW00} R.~M.~Guralnick and J.~S.~Wilson, \textit{The probability of generating a finite soluble group}, Proc.\ London Math.\ Soc.\ \textbf{81} (2000), no.~2, 405--427.
\bibitem[Gus73]{Gus73} W.~H.~Gustafson, \textit{What is the probability that two group elements commute?}, Amer.\ Math.\ Monthly \textbf{80} (1973), 1031--1034.
\bibitem[bln]{bln} T. Burness, A. Lucchini and D. Nemmi, \textit{On the soluble graph of a finite group}, J. Combin. Theory Ser. A \textbf{194} (2023), Paper No. 105708, 39 pp.
\bibitem[dh]{dh} K. Doerk and T. Hawkes, \textit{Finite soluble groups}, De Gruyter Expositions in Mathematics, 4. Walter de Gruyter \& Co., Berlin, 1992.
\bibitem[hup]{hup} B. Huppert, \textit{Endliche Gruppen}, Die Grundlehren der mathematischen Wissenschaften, Band 134. Springer-Verlag, Berlin-New York, 1967.
\bibitem[mqrd]{mqrd} N. Menezes, M. Quick and C. Roney-Dougal, \textit{The probability of generating a finite simple group}, Israel J. Math. \textbf{198}(1) (2013) 371–392.
\bibitem[pf]{pf} P. Flavell, \textit{Finite groups in which every two elements generate a soluble group}, Invent. Math. \textbf{121} (1995) 279-285.
\bibitem[th]{th} J. G. Thompson, \textit{Nonsolvable finite groups all of whose local subgroups are solvable}, Bull. Amer. Math. Soc. \textbf{74} (1968) 383-437.
\end{thebibliography}
\end{document}
